\begin{definition}[Operators generated by tensors with fusion isometries]%
    \label{def:mpdo_bnt_fusion_operator_family}
    \lean{MPOTensor.BNTFusionIsometryFamily.toOperatorFamily}
    \lean{MPOTensor.BNTFusionIsometryFamily.toOperatorFamily_operator}
    \leanok
    \uses{def:mpdo_bnt_fusion_isometry_family, def:mpdo_bnt_label_operator_family, def:mpdo_positive_bnt_label_chi_trace_power_form}
    The tensors $M_\gamma$ generate the periodic operators $O_L(M_\gamma)$.
    The positive diagonal matrices $\chi_{\alpha,\beta,\gamma}$ define the
    structure coefficients
    $c^{(L)}_{\alpha,\beta,\gamma}
        = \tr(\chi_{\alpha,\beta,\gamma}^{L})$.
\end{definition}

\begin{figure}[ht]
    \centering
    % Formula:
    %   O_L(M_\alpha)O_L(M_\beta)
    %     = \sum_\gamma \tr
    %       (\chi_{\alpha,\beta,\gamma}^{L})O_L(M_\gamma).
    % Source verdict: CPSV16 Theorem 4.14(ii)--(iii), source TeX lines
    %   962--991, states the formula algebraically.  MPDO_O_L_A_.png fixes
    %   the periodic-word convention for O_L(M); the source has no dedicated
    %   picture of the complete product law.  The horizontal words below are
    %   the house-style quarter-turn of that source word.
    % Ink-to-index map: at site \ell on the left, the vertical pairleg carries
    %   q_\ell and contracts the bra index of M_\alpha with the ket index of
    %   M_\beta.  Their periodic rows contract a_0=a_L and b_0=b_L
    %   independently; i_\ell and j_\ell remain free.  On the right, the
    %   upper loop is the multiplicity trace
    %   tr(\chi_{\alpha,\beta,\gamma}^{L}) (labelled \chi in the ink), and the
    %   lower loop is O_L(M_\gamma).  The two loops share no index, so each
    %   is a whole picture, and their juxtaposition is the product of a
    %   scalar and an operator.
    %   The sum over \gamma is a sector sum, not a contraction edge.
    % Boundary signatures (virtual-west, virtual-east, physical-up,
    %   physical-down): both sides have semantic signature (0,0,L,L).
    %   Each three-site representative emits (0,0,3,3).
    % The elided columns contract in pairs like the displayed ones, and the
    %   frame would mint that pairleg between the two ellipsis cells as a
    %   stub; the left panel therefore authors its bonds and leaves the
    %   ellipsis column bare.
    \tenkzkernel
    \begin{tenkz}[rows={op,op}, cols=4, west=trace, east=trace, bonds=none]
        \tn[skin=mpo, ports={90:physical:$i_1$, 270:physical}]{M_\alpha} &
        \tn[skin=mpo, ports={90:physical, 270:physical}]{M_\alpha} &
        \tn[skin=dots]{} &
        \tn[skin=mpo, ports={90:physical:$i_L$, 270:physical}]{M_\alpha} \\
        \tn[skin=mpo, ports={90:physical, 270:physical:$j_1$}]{M_\beta} &
        \tn[skin=mpo, ports={90:physical, 270:physical}]{M_\beta} &
        \tn[skin=dots]{} &
        \tn[skin=mpo, ports={90:physical, 270:physical:$j_L$}]{M_\beta}
        \tnwire{(1,1)}{(1,2)} \tnwire{(1,2)}{(1,3)} \tnwire{(1,3)}{(1,4)}
        \tnwire{(2,1)}{(2,2)} \tnwire{(2,2)}{(2,3)} \tnwire{(2,3)}{(2,4)}
        \tnwire{(1,1)}{(2,1)} \tnwire{(1,2)}{(2,2)} \tnwire{(1,4)}{(2,4)}
        \tnmark[form=bracket]{(2,1) .. (2,4)}{$L\text{ sites}$}
    \end{tenkz}
    \quad$=\sum_\gamma$\quad
    \begin{tabular}[c]{@{}c@{}}
        {\begin{tenkz}[cols=4, west=trace, east=trace]
             \tn[skin=box]{\chi} & \tn[skin=box]{\chi} & \tn[skin=dots]{} &
             \tn[skin=box]{\chi}
         \end{tenkz}} \\
        {\begin{tenkz}[cols=4, west=trace, east=trace]
             \tn[skin=mpo,
                 ports={90:physical:$i_1$, 270:physical:$j_1$}]{M_\gamma} &
             \tn[skin=mpo, ports={90:physical, 270:physical}]{M_\gamma} &
             \tn[skin=dots]{} &
             \tn[skin=mpo,
                 ports={90:physical:$i_L$, 270:physical:$j_L$}]{M_\gamma}
             \tnmark[form=bracket]{(1,1) .. (1,4)}{$L\text{ sites}$}
         \end{tenkz}}
    \end{tabular}
    \caption{Closing an $L$-site chain of the local fusion identity. On the
        left, the bra legs of the $M_\alpha$ operator are contracted with the
        ket legs of the $M_\beta$ operator. On the right, the multiplicity loop
        and the $M_\gamma$ loop are disjoint; the former contributes
        $\tr(\chi_{\alpha,\beta,\gamma}^{L})$. Thus the diagram is
        the trace-power product law of
        \cite[Theorem~4.14(ii)--(iii), equations~(87) and~(89),
            lines~962--991]{Cirac2016MPDO_arXiv}.}
    \label{fig:mpdo_fusion_trace_power}
\end{figure}

\begin{theorem}[Fusion identities give the operator product law]%
    \label{thm:mpdo_fusion_same_length_product}
    \lean{MPOTensor.BNTFusionCoisometryFamily.mpo_mul_mpo_eq_sum}
    \lean{MPOTensor.BNTFusionIsometryFamily.mpo_mul_mpo_eq_sum}
    \lean{MPOTensor.BNTFusionIsometryFamily.toOperatorFamily_hasSameLengthProductForm}
    \leanok
    \uses{def:mpdo_bnt_fusion_coisometry_family, def:mpdo_bnt_fusion_isometry_family, def:mpdo_bnt_fusion_operator_family, thm:mpdo_operator_product_mpo, lem:mpdo_word_product_conjugation, lem:mpdo_word_product_block_kronecker, def:mpdo_bnt_label_same_length_product_form}
    Fusion coisometries onto the active sectors with exact reconstruction,
    and also fusion isometries on the full product bond space, give the
    operator product law with structure coefficients determined by the
    diagonal matrices $\chi_{\alpha,\beta,\gamma}$. For every chain
    length $L\ge 1$,
    \begin{align}
        O_L(M_\alpha)O_L(M_\beta)
         & =\sum_{\gamma}
        \tr(\chi_{\alpha,\beta,\gamma}^{L})O_L(M_\gamma).
        \label{eq:rfp_same_length_product}
    \end{align}
    This derives the product law of
    \cite[Theorem~4.14(ii)]{Cirac2016MPDO_arXiv} from the isometry statement
    of \cite[Theorem~4.14(iii)]{Cirac2016MPDO_arXiv}. The active-support
    formulation is the one used in
    \cite[Appendix~C.4, lines~2020--2029]{Cirac2016MPDO_arXiv}.
\end{theorem}
\begin{proof}\leanok
    The product of the operators is generated by the product tensor. Write
    $P^w$ for a nonempty word in that tensor and $D^w$ for the corresponding
    word in the weighted direct sum. For a coisometry onto the active
    sectors, exact reconstruction and
    $U_{\alpha,\beta}U_{\alpha,\beta}^\dagger=1$ give
    $P^w=U_{\alpha,\beta}^\dagger D^w U_{\alpha,\beta}$.
    Cyclicity of trace then gives $\tr(P^w)=\tr(D^w)$.
    For an isometry on the full product bond space,
    $D^w=U_{\alpha,\beta}P^wU_{\alpha,\beta}^\dagger$ and
    $\tr(D^w)=\tr(P^wU_{\alpha,\beta}^\dagger U_{\alpha,\beta})
        =\tr(P^w)$.
    By the fusion identity, the letters of $D$ are block-diagonal
    tensor products. Their word product splits block by block with the
    diagonal factors accumulating into $\chi_{\alpha,\beta,\gamma}^{L}$;
    taking the trace block by block and factor by factor yields the
    coefficients $\tr(\chi_{\alpha,\beta,\gamma}^{L})$ in
    (\ref{eq:rfp_same_length_product}).
\end{proof}

\begin{definition}[Algebra relations from fusion isometries]%
    \label{def:mpdo_bnt_fusion_algebra_clause}
    \lean{MPOTensor.BNTFusionIsometryFamily.toBNTAlgebraClause}
    \leanok
    \uses{def:mpdo_bnt_algebra_clause, def:mpdo_bnt_fusion_isometry_family, def:mpdo_bnt_label_idempotent_coefficient_form, def:mpdo_positive_bnt_label_chi_trace_power_form, thm:mpdo_fusion_same_length_product}
    Suppose the fusion-isometry identity holds with positive diagonal matrices
    $\chi_{\alpha,\beta,\gamma}$. Set
    $c^{(L)}_{\alpha,\beta,\gamma}
        =\tr(\chi_{\alpha,\beta,\gamma}^{L})$,
    and suppose that the vector $m=(m_\alpha)_\alpha$ is an idempotent for
    the multiplication induced by $c^{(1)}$.
    For $L>0$, the operators $O_L(M_\alpha)$ linearly span an algebra.
    The product law and the assumed idempotent condition are, respectively,
    \begin{align}
        O_L(M_\alpha)O_L(M_\beta)
         & =\sum_\gamma
        c^{(L)}_{\alpha,\beta,\gamma}O_L(M_\gamma),
        \notag                  \\
        m_\gamma
         & =\sum_{\alpha,\beta}
        c^{(1)}_{\alpha,\beta,\gamma}m_\alpha m_\beta.
        \notag
    \end{align}
    This construction realizes implication (iii)$\Rightarrow$(ii) of
    \cite[Theorem~4.14]{Cirac2016MPDO_arXiv}; it does not assert either of the
    implications involving the renormalization fixed-point condition.
\end{definition}

\begin{definition}[Active-sector fusion for a vertical canonical decomposition]%
    \label{def:mpdo_bnt_fusion_tensor_clause}
    \lean{MPOTensor.BNTFusionTensorClause}
    \lean{MPOTensor.HasBNTFusionTensorClause}
    \leanok
    \uses{def:vertical_cf_mpdo, def:mpdo_bnt_fusion_coisometry_family, def:mpdo_bnt_label_idempotent_coefficient_form}
    Choose a vertical canonical decomposition, with nonempty positive
    diagonal matrices $\mu_\alpha$ and a BNT $\{M_\alpha\}_\alpha$:
    \begin{align}
        U\widetilde M U^\dagger
         & =\bigoplus_\alpha \mu_\alpha\otimes M_\alpha.
        \notag
    \end{align}
    Here $U$ is a coisometry onto the retained sectors, and its adjoint
    reconstructs $\widetilde M$, including any common zero-sector complement.
    Active-sector fusion for this decomposition consists of positive diagonal
    matrices $\chi_{\alpha,\beta,\gamma}$ and coisometries
    $U_{\alpha,\beta}$ onto the active product sectors, satisfying
    \begin{align}
        U_{\alpha,\beta}(M_\alpha M_\beta)^{ij}
        U_{\alpha,\beta}^\dagger
         & =\bigoplus_\gamma
        \chi_{\alpha,\beta,\gamma}\otimes M_\gamma^{ij}.
        \label{eq:rfp_active_fusion}
    \end{align}
    They obey $U_{\alpha,\beta}U_{\alpha,\beta}^\dagger=1$, and the active
    direct sum reconstructs the product tensor:
    \begin{align}
        (M_\alpha M_\beta)^{ij}
         & =U_{\alpha,\beta}^\dagger
        \left(\bigoplus_\gamma
        \chi_{\alpha,\beta,\gamma}\otimes M_\gamma^{ij}\right)
        U_{\alpha,\beta}.
        \notag
    \end{align}
    The vector with entries $m_\alpha=\tr(\mu_\alpha)$ is an idempotent
    for the multiplication induced by
    $c^{(1)}_{\alpha,\beta,\gamma}=\tr(\chi_{\alpha,\beta,\gamma})$.
    This is the active-support formulation of statement (iii) of
    \cite[Theorem~4.14, lines~986--993]{Cirac2016MPDO_arXiv}, for the
    chosen decomposition of
    \cite[Proposition~4.13, lines~943--951]{Cirac2016MPDO_arXiv}.
\end{definition}

\begin{theorem}[Idempotent relation from paired decompositions and given fusion
        coisometries]%
    \label{thm:mpdo_paired_vertical_decompositions_fusion}
    \lean{MPOTensor.HasBNTFusionTensorClause.%
        of_verticalDecompositions_of_unitaryBlockEquiv}
    \leanok
    \uses{def:mpdo_bnt_fusion_tensor_clause}
    Let $\mathcal D_1$ and $\mathcal D_2$ be positive vertical canonical
    decompositions of an MPO tensor and its two-site blocking. Suppose their
    normal sectors are identified by a bijection $\sigma$ and unitary
    conjugacies with the corresponding multiplicity-trace ratios. Suppose also
    that positive diagonal matrices $\chi_{\alpha,\beta,\gamma}$ and
    coisometries $U_{\alpha,\beta}$ are given, with
    $U_{\alpha,\beta}U_{\alpha,\beta}^{\dagger}=1$, and satisfy both
    \begin{align}
        U_{\alpha,\beta}(M_\alpha M_\beta)^{ij}
        U_{\alpha,\beta}^{\dagger}
         & =\bigoplus_\gamma
        \chi_{\alpha,\beta,\gamma}\otimes M_\gamma^{ij},
        \notag                         \\
        (M_\alpha M_\beta)^{ij}
         & =U_{\alpha,\beta}^{\dagger}
        \left(\bigoplus_\gamma
        \chi_{\alpha,\beta,\gamma}\otimes M_\gamma^{ij}\right)
        U_{\alpha,\beta}.
        \notag
    \end{align}
    Then the simultaneous blocked representations imply that the vector
    with entries $m_\alpha=\tr(\mu_\alpha)$ is an idempotent for the
    multiplication induced by
    $c^{(1)}_{\alpha,\beta,\gamma}=\tr(\chi_{\alpha,\beta,\gamma})$.
    Thus the chosen decomposition satisfies the active-sector fusion and
    idempotent relations. With the stated comparison and fusion hypotheses,
    this is the final step of the construction in
    \cite[Appendix~C.4, lines~2001--2046]{Cirac2016MPDO_arXiv}.
\end{theorem}
\begin{proof}\leanok
    \uses{thm:mpdo_blocked_vertical_operator_representations, thm:mpdo_blocked_representations_idempotent_trace_scalars}
    The sector correspondence expresses the blocked closed-chain operator in
    the one-site normal-tensor basis. Multiplication of the one-site vertical
    decomposition gives a second expression in the same basis. At every
    sufficiently large length, linear independence equates their
    coefficients. Equality of the resulting positive power sums identifies
    the multiplicity multisets, and summing their entries gives
    $m_\gamma=\sum_{\alpha,\beta}
        \tr(\chi_{\alpha,\beta,\gamma})m_\alpha m_\beta$.
\end{proof}

\begin{theorem}[Renormalization fixed points in CF have active-sector fusion]%
    \label{thm:mpdo_literal_cpsv_rfp_to_bnt_fusion_tensor}
    \lean{MPOTensor.HasBNTFusionTensorClause.of_isRFPViaTS}
    \leanok
    \uses{def:mpdo, def:cpsv_canonical_form, def:is_rfp_via_ts, def:mpdo_bnt_fusion_tensor_clause}
    Let $M$ generate MPDOs $\rho^{(N)}(M)$, and suppose that its
    doubled-physical-index MPS tensor is in CF, with normalized NT blocks.
    Suppose there are trace-preserving completely positive maps (tpCPM)
    $S,T$ on the original physical indices such that
    $S[M_2(X)]=M_1(X)$ and $T[M_1(X)]=M_2(X)$ for every $X$.
    Then a vertical canonical decomposition of $M$ has active-sector
    fusion. Its positive diagonal fusion matrices satisfy
    \begin{align}
        m_\gamma
         & =\sum_{\alpha,\beta}
        \tr(\chi_{\alpha,\beta,\gamma})m_\alpha m_\beta.
        \notag
    \end{align}
    This is the corrected active-support formulation of implication
    (i)$\Rightarrow$(iii) in
    \cite[Theorem~4.14, lines~972--993]{Cirac2016MPDO_arXiv}, rather than the
    unqualified printed statement.

    For each fixed product pair, only its active normal corners occur;
    an absent label has zero fusion multiplicity. See
    \cite{gap:cpsv16_bnt_uniqueness_zero_coefficient}.

    A fixed product pair may have an empty active family, so no
    unsupported normal sector is added. See
    \cite{gap:cpsv16_figure11_per_pair_support}.

    Each fusion map is a coisometry onto the active direct sum, and its
    adjoint reconstructs the whole product tensor, including any discarded
    common zero corner. This corrects the orientation of the printed map.
    See \cite{gap:cpsv16_figure11_fusion_coisometry}.
\end{theorem}

\begin{proof}\leanok
    \uses{thm:mpdo_literal_cpsv_vertical_decomposition, thm:mpdo_literal_cpsv_vertical_decomposition_block_two, thm:mpdo_transported_vertical_sector_coefficient_comparison, thm:mpdo_literal_cpsv_vertical_product_positive_fusion_decomposition, thm:mpdo_paired_vertical_decompositions_fusion}
    Choose the one-site and two-site vertical decompositions obtained from
    CF. Transporting $T$ and $S$ through them identifies their normal
    sectors by unitary conjugacies scaled by multiplicity-trace ratios.
    The positive-fusion decomposition gives the positive diagonal matrices,
    coisometries, and exact reconstruction on every active product sector.
    The paired-decomposition theorem then gives the idempotent relation for
    $m_\alpha=\tr(\mu_\alpha)$.
\end{proof}

\begin{theorem}[Horizontal-form specialization of active-sector fusion]%
    \label{thm:mpdo_horizontal_rfp_to_bnt_fusion_tensor}
    \lean{MPOTensor.HasBNTFusionTensorClause.of_isRFPViaTS_of_horizontalCF}
    \leanok
    \uses{def:is_rfp_via_ts, def:horizontal_cf_mpdo, def:mpdo_bnt_fusion_tensor_clause}
    Let $M$ generate MPDOs and satisfy the normalized BNT-refined
    horizontal-form hypothesis. Suppose there are tpCPM $S,T$ on the
    original physical indices with $S[M_2(X)]=M_1(X)$ and
    $T[M_1(X)]=M_2(X)$ for every $X$. Then a vertical canonical
    decomposition of $M$ has active-sector fusion. Its positive diagonal
    fusion matrices make the vector $m=(\tr(\mu_\alpha))_\alpha$ an
    idempotent for the multiplication induced by $c^{(1)}$.

    The normalized BNT-refined horizontal hypothesis is stronger than
    CF with normalized NT blocks.
    See \cite{gap:cpgsv17_vertical_cf_grouping}.
\end{theorem}
\begin{proof}\leanok
    \uses{thm:vertical_cf_of_horizontal_cf, thm:mpdo_transported_vertical_sector_coefficient_comparison, thm:mpdo_transported_vertical_product_positive_fusion_decomposition, thm:mpdo_paired_vertical_decompositions_fusion}
    Choose the one-site and two-site vertical canonical decompositions supplied
    by the normalized BNT-refined horizontal form. The renormalization maps
    identify their normal sectors and give positive diagonal fusion matrices
    and fusion coisometries.
    Apply the paired-decomposition theorem to obtain the idempotent relation
    for $m_\alpha=\tr(\mu_\alpha)$.
\end{proof}

\begin{theorem}[Active-sector fusion implies the algebra relations]%
    \label{thm:mpdo_bnt_fusion_tensor_to_algebra_tensor}
    \lean{MPOTensor.BNTFusionTensorClause.toBNTAlgebraTensorClause}
    \lean{MPOTensor.HasBNTFusionTensorClause.to_hasBNTAlgebraTensorClause}
    \leanok
    \uses{def:mpdo_bnt_algebra_tensor_clause, def:mpdo_bnt_fusion_algebra_clause, def:mpdo_bnt_fusion_tensor_clause}
    Suppose a chosen vertical canonical decomposition satisfies the
    active-sector fusion and idempotent relations.
    Then its operators $O_L(M_\alpha)$ satisfy the algebra relations.
    For every positive chain length $L$, the structure coefficients are
    $c^{(L)}_{\alpha,\beta,\gamma}
        =\tr(\chi_{\alpha,\beta,\gamma}^{L})$,
    determined by the same diagonal matrices. Thus statement (iii) implies
    statement (ii) of
    \cite[Theorem~4.14, lines~972--993]{Cirac2016MPDO_arXiv}, with the
    active-support interpretation of (iii).
    No assertion concerning the renormalization fixed-point condition is
    made here.
\end{theorem}
\begin{proof}\leanok
    Iterating the fusion identity along $L$ sites and taking the closed trace
    gives
    $O_L(M_\alpha)O_L(M_\beta)
        =\sum_\gamma
        \tr(\chi_{\alpha,\beta,\gamma}^{L})O_L(M_\gamma)$.
    Writing $m_\alpha=\tr(\mu_\alpha)$, the assumed
    idempotent identity is
    $m_\gamma=\sum_{\alpha,\beta}
        \tr(\chi_{\alpha,\beta,\gamma})m_\alpha m_\beta$.
    These are the product and idempotent relations for the chosen
    vertical decomposition, with the same diagonal matrices.
\end{proof}

% CPSV16 Theorem 4.14 (`thm:IV.13`), lines 972--993 and Appendix C.4.
% Gaps: docs/paper-gaps/cpsv16_figure11_per_pair_support.tex and
% docs/paper-gaps/cpsv16_figure11_fusion_coisometry.tex.
\begin{theorem}[Printed algebra and fusion characterization]
    \label{thm:cpsv_theorem414_printed_status}
    \notready
    \uses{def:mpdo, def:cpsv_canonical_form, def:is_rfp_via_ts, def:mpdo_bnt_algebra_tensor_clause, def:mpdo_bnt_fusion_tensor_clause, thm:mpdo_cpsv_bnt_rfp_equivalence}
    For a tensor in CF generating MPDOs,
    \cite[Theorem~4.14]{Cirac2016MPDO_arXiv} states that the RFP
    condition, the algebra and idempotent relations for a vertical BNT,
    and the fusion-isometry and idempotent relations are equivalent.

    The RFP condition is equivalent to the algebra and idempotent
    relations. For the fusion statement, an absent product sector must have
    multiplicity zero, and the map onto the retained sectors is a coisometry
    with exact reconstruction. Under this active-support interpretation,
    the three conditions are equivalent by
    Theorem~\ref{thm:mpdo_cpsv_bnt_rfp_equivalence}. Without these
    qualifications, the displayed fusion statement is not well-defined on
    empty product sectors and has the wrong isometry orientation.
\end{theorem}

\begin{theorem}[BNT characterizations of MPDO renormalization fixed points]
    \label{thm:mpdo_cpsv_bnt_rfp_equivalence}
    \lean{MPOTensor.isRFPViaTS_iff_hasBNTAlgebraTensorClause}
    \lean{MPOTensor.isRFPViaTS_iff_hasBNTFusionTensorClause}
    \leanok
    \uses{def:mpdo, def:cpsv_canonical_form, def:is_rfp_via_ts, def:mpdo_bnt_algebra_tensor_clause, def:mpdo_bnt_fusion_tensor_clause}
    Let $M$ generate MPDOs, and suppose its doubled-physical-index tensor is
    in CF with normalized NT blocks. Consider the following conditions:
    \begin{enumerate}
        \item[(i)] $M$ is an RFP: there are tpCPM $S,T$ with
              $S[M_2(X)]=M_1(X)$ and $T[M_1(X)]=M_2(X)$ for every $X$.
        \item[(ii)] A vertical canonical decomposition of $M$ has positive
              diagonal matrices $\chi_{\alpha,\beta,\gamma}$ such that
              the $O_L(M_\alpha)$ satisfy the product law with
              $c^{(L)}_{\alpha,\beta,\gamma}
                  =\tr(\chi_{\alpha,\beta,\gamma}^{L})$ for every $L>0$,
              and the vector with entries $m_\alpha=\tr(\mu_\alpha)$ is an
              idempotent for the multiplication induced by $c^{(1)}$.
        \item[(iii$_{\mathrm{act}}$)] A vertical canonical decomposition of
              $M$ has active-sector fusion coisometries with exact
              reconstruction and positive diagonal matrices
              $\chi_{\alpha,\beta,\gamma}$, and the vector with entries
              $m_\alpha=\tr(\mu_\alpha)$ is an idempotent for the
              multiplication induced by $c^{(1)}$.
    \end{enumerate}
    Then
    \begin{align}
        \textnormal{(i)} & \iff\textnormal{(ii)}, \notag                       \\
        \textnormal{(i)} & \iff\textnormal{(iii}_{\mathrm{act}}\textnormal{)}.
        \label{eq:mpdo_cpsv_bnt_rfp_equivalence}
    \end{align}
    The first equivalence is the equivalence of (i) and (ii) in
    \cite[Theorem~4.14, lines~972--985]{Cirac2016MPDO_arXiv}. The second is
    the corrected active-support form corresponding to the printed
    conditions (i) and (iii), rather than an assertion of the unrestricted
    printed condition (iii).

    Only nonzero product corners occur in (iii$_{\mathrm{act}}$);
    a sector absent from a fixed product pair has zero fusion multiplicity.
    This restriction is recorded in
    \cite{gap:cpsv16_bnt_uniqueness_zero_coefficient}.

    A fixed pair may have empty active support, and the fusion map onto
    the retained sectors is a coisometry. These corrections are recorded in
    \cite{gap:cpsv16_figure11_per_pair_support} and
    \cite{gap:cpsv16_figure11_fusion_coisometry}.

    The implication (ii)$\Rightarrow$(i) uses the comparison of mixed
    one-site and two-site prefixes recorded in
    \cite{gap:cpsv16_two_site_sector_unitary_gauge_gap}.

    The physical maps in (ii)$\Rightarrow$(i) are completed on the
    discarded complements as recorded in
    \cite{gap:cpgsv17_vertical_isometry_zero_sector}.
\end{theorem}

\begin{proof}\leanok
    \uses{thm:mpdo_literal_cpsv_rfp_to_bnt_fusion_tensor, thm:mpdo_bnt_fusion_tensor_to_algebra_tensor, thm:mpdo_bnt_algebra_tensor_to_rfp}
    An RFP has active-sector fusion and the idempotent relation by
    Theorem~\ref{thm:mpdo_literal_cpsv_rfp_to_bnt_fusion_tensor}.
    These imply the algebra and idempotent relations by
    Theorem~\ref{thm:mpdo_bnt_fusion_tensor_to_algebra_tensor}.
    Finally, the algebra and idempotent relations imply that $M$ is an RFP by
    Theorem~\ref{thm:mpdo_bnt_algebra_tensor_to_rfp}. These three
    implications prove both equivalences in
    (\ref{eq:mpdo_cpsv_bnt_rfp_equivalence}).
\end{proof}

\begin{figure}[ht]
    \centering
    % Writing Q_N for the N-site operator and r=N-1, the recursion is
    %   Q_N = \bigoplus_{\substack{\alpha_1,\ldots,\alpha_{r+1}\\
    %                    \beta_1,\ldots,\beta_r\\k_1,\ldots,k_r}}
    %         \left(\prod_{s=1}^{r}
    %           \chi_{\alpha_s,\beta_s,\alpha_{s+1},k_s}\right)
    %         \tr(M_{\alpha_{r+1}}).
    % Source verdict: CPSV16 lines 999--1008 and MPDO_Structure.png give this
    %   recursive coefficient graph.  The image is the source for the open
    %   copy rails, repeated-step ellipsis, and terminal horizontal trace.
    % Ink-to-index map: each X_s casts four east legs, one per index
    %   k_s, alpha_s, beta_s, gamma_s, each ending in a dot on an open copy
    %   rail.  For consecutive steps, gamma_s and alpha_{s+1} are two
    %   incidences on the same rail, which remains open above and below; the
    %   rail is a diagonal copy register and does not sum that label.  A leg
    %   passing a rail without a dot does not tap it.  The
    %   ellipsis repeats precisely this routing.  On the terminal rail,
    %   gamma_r meets both X_r and the selector dot feeding M_{gamma_r}.
    %   The two operator ports of M_{gamma_r} leave its lower corners and are
    %   contracted with one another below it, while its two bond ports remain
    %   open above and below.  Thus the terminal
    %   matrix acts on the label-dependent bond space of M_{gamma_r}.
    % Boundary signature: the r coefficient boxes have 3r+1 unique open copy
    %   rails, together with the open terminal bond pair.  The bond dimension
    %   depends on the terminal label gamma_r.
    % Projector caveat: lines 1010--1015 require length-independent structure
    %   coefficients and projector terminal traces before Q is a projector;
    %   the contraction displayed here asserts only the operator Q_N.
    \ensuremath{Q={}}
    \tenkzkernel
    \begin{tenkz}[rows={wire,wire,wire,wire,wire,wire}, cols=19, bonds=none,
            metrics=compact]
        % Nineteen columns at the base metric outrun the compiled page and
        % lose the terminal trace; the compact profile keeps the whole
        % recursion on one line.
        \tn[at=(5,1), skin=box, name=X1,
            ports={45:physical, 0:physical, -20:physical, -45:physical}]{X}
        \tn[at=(4,5), skin=box, name=X2,
            ports={45:physical, 0:physical, -20:physical, -45:physical}]{X}
        \tn[at=(3,9), skin=box, name=X3,
            ports={45:physical, 0:physical, -20:physical, -45:physical}]{X}
        \tn[at=(2,14), skin=box, name=X4,
            ports={45:physical, 0:physical, -20:physical, -45:physical}]{X}
        \tn[at=(4,2), name=k1, ports={180:physical, 90:physical, 270:physical}]{}
        \tn[at=(5,3), name=a1, ports={180:physical, 90:physical, 270:physical}]{}
        \tn[at=(6,4), name=b1, ports={180:physical, 90:physical, 270:physical}]{}
        \tn[at=(3,6), name=k2, ports={180:physical, 90:physical, 270:physical}]{}
        \tn[at=(4,7), name=a2, ports={180:physical, 90:physical, 270:physical}]{}
        \tn[at=(5,8), name=b2, ports={180:physical, 90:physical, 270:physical}]{}
        \tn[at=(6,7), name=g1, ports={180:physical, 90:physical, 270:physical}]{}
        \tn[at=(2,10), name=k3, ports={180:physical, 90:physical, 270:physical}]{}
        \tn[at=(3,11), name=a3, ports={180:physical, 90:physical, 270:physical}]{}
        \tn[at=(4,12), name=b3, ports={180:physical, 90:physical, 270:physical}]{}
        \tn[at=(5,11), name=g2, ports={180:physical, 90:physical, 270:physical}]{}
        \tn[at=(4,13), name=g3, ports={180:physical, 90:physical, 270:physical}]{}
        \tn[at=(1,15), name=k4, ports={180:physical, 90:physical, 270:physical}]{}
        \tn[at=(2,16), name=a4, ports={180:physical, 90:physical, 270:physical}]{}
        \tn[at=(3,17), name=b4, ports={180:physical, 90:physical, 270:physical}]{}
        \tn[at=(3,18), name=g4, ports={180:physical, 90:physical, 270:physical}]{}
        \tn[at=(4,18), name=sel, ports={90:physical, 270:physical, 0:physical}]{}
        \tn[at={1 e of sel}, skin=box, name=M,
            ports={180:physical, 90:virtual, 270:virtual,
                    225:physical, 315:physical}]{M_{\gamma_r}}
        \tn[at=(5,19), skin=none, name=tr, ports={135:physical, 45:physical}]{}
        \tn[at={midway (2,12) and (3,13)}, skin=dots, name=dd]{}
        \tnwire{X1.45}{k1.180}
        \tnwire{X1.0}{a1.180}
        \tnwire{X1.-45}{b1.180}
        \tnwire{X1.-20}{g1.180}
        \tnwire{X2.45}{k2.180}
        \tnwire{X2.0}{a2.180}
        \tnwire{X2.-45}{b2.180}
        \tnwire{X2.-20}{g2.180}
        \tnwire{X3.45}{k3.180}
        \tnwire{X3.0}{a3.180}
        \tnwire{X3.-45}{b3.180}
        \tnwire{X3.-20}{g3.180}
        \tnwire{X4.45}{k4.180}
        \tnwire{X4.0}{a4.180}
        \tnwire{X4.-45}{b4.180}
        \tnwire{X4.-20}{g4.180}
        \tnwire{k1.90}{open n}
        \tnwire{k1.270}{open s}
        \tnwire{a1.90}{open n}
        \tnwire{a1.270}{open s}
        \tnwire{b1.90}{open n}
        \tnwire{b1.270}{open s}
        \tnwire{k2.90}{open n}
        \tnwire{k2.270}{open s}
        \tnwire{a2.90}{open n}
        \tnwire{a2.270}{g1.90}
        \tnwire{g1.270}{open s}
        \tnwire{b2.90}{open n}
        \tnwire{b2.270}{open s}
        \tnwire{k3.90}{open n}
        \tnwire{k3.270}{open s}
        \tnwire{a3.90}{open n}
        \tnwire{a3.270}{g2.90}
        \tnwire{g2.270}{open s}
        \tnwire{b3.90}{open n}
        \tnwire{b3.270}{open s}
        \tnwire{g3.90}{open n}
        \tnwire{g3.270}{open s}
        \tnwire{k4.90}{open n}
        \tnwire{k4.270}{open s}
        \tnwire{a4.90}{open n}
        \tnwire{a4.270}{open s}
        \tnwire{b4.90}{open n}
        \tnwire{b4.270}{open s}
        \tnwire{g4.90}{open n}
        \tnwire{g4.270}{sel.90}
        \tnwire{sel.270}{open s}
        \tnwire{M.180}{sel.0}
        \tnwire{M.90}{open n}
        \tnwire{M.270}{open s}
        \tnwire{M.225}{tr.135}
        \tnwire{M.315}{tr.45}
    \end{tenkz}
    \caption{The recursive contraction defining the structure operator $Q$,
        with terminal factor $\tr(M_{\gamma_r})$
        \cite[lines~999--1008]{Cirac2016MPDO_arXiv}.}
    \label{fig:mpdo_recursive_structure_operator}
\end{figure}

The associativity of the algebra product in
\cite[Equation~(87)]{Cirac2016MPDO_arXiv}
constrains the diagonal matrices $\chi_{\alpha,\beta,\gamma}$ and, through the
isometry statement, induces a compatibility equation between the fusion
isometries of the two ways of bracketing a triple product,
$(M_\alpha M_\beta) M_\gamma$ and $M_\alpha (M_\beta M_\gamma)$
\cite[lines~995--999]{Cirac2016MPDO_arXiv}. The source relates solutions
constructed from the $F$-symbols of unitary fusion categories to the
boundary theories studied in
\cite{Bultinck2017Anyons}. The following results concern fusion isometries
on the full product bond spaces. They give the tensor identities for the
two bracketings and compare the resulting full direct sums.
Restriction of that comparison to a fixed final label is a separate step.

\begin{figure}[ht]
    \centering
    % The parenthesization is the tree: delta and eta label the respective
    % internal edges, while epsilon labels the common root edge.
    \begin{align}
        \tntree{((\alpha\,\beta)_\delta\,\gamma)_\varepsilon}
         & \xrightarrow{F_{\varepsilon}^{\alpha\beta\gamma}}
        \tntree{(\alpha\,(\beta\,\gamma)_\eta)_\varepsilon}.
        \notag
    \end{align}
    \caption{The two ways of fusing three labels to a fixed final label
        $\varepsilon$, using the printed-$F$ row/column convention of
        \cite[Section~3.4]{Bultinck2017Anyons}. The solid fusion trees are
        the left- and right-bracketed maps constructed below. The arrow is
        the change of basis from left- to right-bracketed multiplicity
        coordinates. The results below first construct this map and then
        prove its two-sided unitarity under the
        simultaneous-selector and positive-bond-dimension hypotheses.}
    \label{fig:mpdo_fixed_final_fusion_bracketings}
\end{figure}

\begin{definition}[Iterated fusion isometry of a triple product, left
        bracketing]%
    \label{def:mpdo_left_fusion_isometry}
    \lean{MPOTensor.BNTFusionIsometryFamily.leftFusionIsometry}
    \leanok
    \uses{def:mpdo_bnt_fusion_isometry_family}
    For labels $\alpha, \beta, \gamma$, the \emph{left iterated fusion
        isometry} $U^{\mathrm L}_{\alpha,\beta,\gamma}$ is obtained by applying
    $U_{\alpha,\beta}$ (tensored with the identity on the bond space of
    $M_\gamma$) to fuse the first two tensors, then applying $U_{\delta,\gamma}$,
    for each label $\delta$ appearing in the resulting direct sum, to fuse
    the outcome with the third tensor. Its codomain is indexed by pairs of
    labels $(\delta, \varepsilon)$, the multiplicity indices of
    $\chi_{\alpha,\beta,\delta}$ and $\chi_{\delta,\gamma,\varepsilon}$,
    and the bond index of $M_\varepsilon$.
\end{definition}

\begin{theorem}[The left iterated fusion isometry is an isometry]%
    \label{thm:mpdo_left_fusion_isometry_isometry}
    \lean{MPOTensor.BNTFusionIsometryFamily.leftFusionIsometry_isometry}
    \leanok
    \uses{def:mpdo_left_fusion_isometry}
    The left iterated fusion isometry satisfies
    $(U^{\mathrm L}_{\alpha,\beta,\gamma})^\dagger
        U^{\mathrm L}_{\alpha,\beta,\gamma}=1$.
\end{theorem}
\begin{proof}\leanok
    The composite of two isometries is an isometry.
\end{proof}

\begin{theorem}[The left iterated fusion isometry conjugates a triple product
        onto a doubly weighted direct sum]%
    \label{thm:mpdo_left_fusion_apply}
    \lean{MPOTensor.BNTFusionIsometryFamily.leftFusion_apply}
    \leanok
    \uses{def:mpdo_left_fusion_isometry, def:mpdo_operator_product, lem:matrix_block_conjugation_restriction}
    Site by site,
    \begin{align}
         & U^{\mathrm L}_{\alpha,\beta,\gamma}
        ((M_\alpha M_\beta)M_\gamma)^{ij}
        (U^{\mathrm L}_{\alpha,\beta,\gamma})^\dagger
        \notag                                 \\
         & \qquad
        =\bigoplus_{\delta}\bigoplus_{\varepsilon}
        \chi_{\alpha,\beta,\delta}
        \otimes\bigl(\chi_{\delta,\gamma,\varepsilon}
        \otimes M_\varepsilon^{ij}\bigr).
        \notag
    \end{align}
\end{theorem}
\begin{proof}\leanok
    \uses{lem:matrix_block_conjugation_restriction}
    Apply the fusion identity for $U_{\alpha,\beta}$ to the first two
    tensors, summed against the physical index contracting with the third
    tensor. Regroup the result over $\delta$ and use blockwise conjugation to
    obtain $\chi_{\alpha,\beta,\delta}$ tensored with a letter of
    $M_\delta M_\gamma$. The fusion identity for $U_{\delta,\gamma}$ gives the
    second direct sum.
\end{proof}

\begin{definition}[Iterated fusion isometry of a triple product, right
        bracketing]%
    \label{def:mpdo_right_fusion_isometry}
    \lean{MPOTensor.BNTFusionIsometryFamily.rightFusionIsometry}
    \leanok
    \uses{def:mpdo_bnt_fusion_isometry_family}
    For labels $\alpha, \beta, \gamma$, the \emph{right iterated fusion
        isometry} $U^{\mathrm R}_{\alpha,\beta,\gamma}$ is obtained by applying
    $U_{\beta,\gamma}$, tensored with the identity on the bond space of
    $M_\alpha$, and then applying $U_{\alpha,\delta}$ for every intermediate
    label $\delta$. Its codomain is indexed by the intermediate and final
    labels $(\delta,\varepsilon)$, the multiplicity indices of
    $\chi_{\beta,\gamma,\delta}$ and $\chi_{\alpha,\delta,\varepsilon}$,
    and the bond index of $M_\varepsilon$.
    This is the right-bracketed composite in the associativity equation of
    \cite[Section~``Associativity and the pentagon equation'']{Bultinck2017Anyons},
    before the change of basis by an $F$-matrix.
\end{definition}

\begin{theorem}[The right iterated fusion isometry is an isometry]%
    \label{thm:mpdo_right_fusion_isometry_isometry}
    \lean{MPOTensor.BNTFusionIsometryFamily.rightFusionIsometry_isometry}
    \leanok
    \uses{def:mpdo_right_fusion_isometry}
    The right iterated fusion isometry satisfies
    $(U^{\mathrm R}_{\alpha,\beta,\gamma})^\dagger
        U^{\mathrm R}_{\alpha,\beta,\gamma}=1$.
\end{theorem}
\begin{proof}\leanok
    The two successive fusion maps are isometries on each intermediate
    sector, so their block-diagonal composite is an isometry.
\end{proof}

\begin{theorem}[The right iterated fusion isometry conjugates a triple product
        onto a doubly weighted direct sum]%
    \label{thm:mpdo_right_fusion_apply}
    \lean{MPOTensor.BNTFusionIsometryFamily.rightFusion_apply}
    \leanok
    \uses{def:mpdo_right_fusion_isometry, def:mpdo_operator_product, lem:matrix_block_conjugation_restriction}
    Site by site,
    \begin{align}
         & U^{\mathrm R}_{\alpha,\beta,\gamma}
        (M_\alpha(M_\beta M_\gamma))^{ij}
        (U^{\mathrm R}_{\alpha,\beta,\gamma})^\dagger
        \notag                                 \\
         & \qquad
        =\bigoplus_{\delta}\bigoplus_{\varepsilon}
        \chi_{\beta,\gamma,\delta}
        \otimes\bigl(\chi_{\alpha,\delta,\varepsilon}
        \otimes M_\varepsilon^{ij}\bigr).
        \notag
    \end{align}
\end{theorem}
\begin{proof}\leanok
    \uses{lem:matrix_block_conjugation_restriction}
    Apply the fusion identity for $U_{\beta,\gamma}$ to $M_\beta M_\gamma$.
    In the summand indexed by $\delta$, exchange the first multiplicity
    factor with the bond space of $M_\alpha$. Blockwise conjugation regroups
    the contraction as
    $\chi_{\beta,\gamma,\delta}\otimes
        (M_\alpha M_\delta)^{ij}$.
    The fusion identity for $U_{\alpha,\delta}$ gives the sum over
    $\varepsilon$.
\end{proof}

\begin{definition}[Full comparison of the two triple-fusion orders]%
    \label{def:mpdo_triple_fusion_comparison}
    \lean{MPOTensor.BNTFusionIsometryFamily.tripleFusionComparison}
    \leanok
    \uses{def:mpdo_left_fusion_isometry, def:mpdo_right_fusion_isometry, def:mpdo_operator_product_associator}
    The canonical comparison from the right-associated full direct sum to the
    left-associated full direct sum is
    \begin{align}
        C_{\alpha,\beta,\gamma}
         & =U^{\mathrm L}_{\alpha,\beta,\gamma}
        A_{D_\alpha,D_\beta,D_\gamma}
        (U^{\mathrm R}_{\alpha,\beta,\gamma})^\dagger.
        \label{eq:rfp_triple_comparison}
    \end{align}
    Here $A_{D_\alpha,D_\beta,D_\gamma}$ is the permutation matrix identifying
    the two bracketings of the product bond space.
    Both intermediate labels and the final label remain in the direct sums.
    This definition does not select a fixed final label and makes no
    invertibility assertion. Its orientation is from the right-associated
    sum to the left-associated sum; consequently it has the orientation of
    the adjoint of the $F$-matrix in the row/column convention used here for
    \cite[Section~3.4]{Bultinck2017Anyons}.
\end{definition}

\begin{lemma}[Left range projection of the full triple-fusion comparison]%
    \label{lem:mpdo_triple_fusion_comparison_left_range}
    \lean{MPOTensor.BNTFusionIsometryFamily.%
        tripleFusionComparison_mul_conjTranspose}
    \leanok
    \uses{def:mpdo_triple_fusion_comparison}
    The product of the full comparison with its adjoint is the range
    projection of the left iterated fusion isometry:
    \begin{align}
        C_{\alpha,\beta,\gamma}C_{\alpha,\beta,\gamma}^{\dagger}
         & =U^{\mathrm L}_{\alpha,\beta,\gamma}
        (U^{\mathrm L}_{\alpha,\beta,\gamma})^{\dagger}.
        \notag
    \end{align}
    The right side is not asserted to be the identity. Surjectivity of the
    left iterated fusion map requires completeness of the fusion
    decomposition, corresponding to the invertible total fusion matrix in
    \cite[lines~181--191]{Bultinck2017Anyons}.
\end{lemma}
\begin{proof}\leanok
    \uses{thm:mpdo_right_fusion_isometry_isometry, def:mpdo_operator_product_associator}
    Substitute (\ref{eq:rfp_triple_comparison}). The isometry identity
    $(U^{\mathrm R})^\dagger U^{\mathrm R}=1$ and the two-sided inverse
    identities for the permutation matrix $A$ leave
    $U^{\mathrm L}(U^{\mathrm L})^\dagger$.
\end{proof}

\begin{lemma}[Right range projection of the full triple-fusion comparison]%
    \label{lem:mpdo_triple_fusion_comparison_right_range}
    \lean{MPOTensor.BNTFusionIsometryFamily.%
        conjTranspose_mul_tripleFusionComparison}
    \leanok
    \uses{def:mpdo_triple_fusion_comparison}
    The product in the opposite order is the range projection of the right
    iterated fusion isometry:
    \begin{align}
        C_{\alpha,\beta,\gamma}^{\dagger}C_{\alpha,\beta,\gamma}
         & =U^{\mathrm R}_{\alpha,\beta,\gamma}
        (U^{\mathrm R}_{\alpha,\beta,\gamma})^{\dagger}.
        \notag
    \end{align}
    The right side is not asserted to be the identity. Surjectivity of the
    right iterated fusion map likewise requires completeness of the fusion
    decomposition.
\end{lemma}
\begin{proof}\leanok
    \uses{thm:mpdo_left_fusion_isometry_isometry, def:mpdo_operator_product_associator}
    Substitute (\ref{eq:rfp_triple_comparison}). The isometry identity
    $(U^{\mathrm L})^\dagger U^{\mathrm L}=1$ and the two-sided inverse
    identities for $A$ leave
    $U^{\mathrm R}(U^{\mathrm R})^\dagger$.
\end{proof}

\begin{theorem}[The full triple-fusion comparison intertwines the unweighted
        direct sums]%
    \label{thm:mpdo_triple_fusion_comparison_intertwines}
    \lean{MPOTensor.BNTFusionIsometryFamily.%
        tripleFusionComparison_intertwines_of_lengthIndependent}
    \leanok
    \uses{def:mpdo_triple_fusion_comparison, thm:mpdo_operator_product_associativity, thm:mpdo_left_fusion_apply, thm:mpdo_right_fusion_apply, thm:mpdo_bnt_length_independent_zipper}
    Suppose the structure coefficients
    $c^{(L)}_{\alpha,\beta,\gamma}
        =\tr(\chi_{\alpha,\beta,\gamma}^{L})$, with positive diagonal
    $\chi_{\alpha,\beta,\gamma}$, are independent of the positive chain
    length. Write $r_{\alpha,\beta}^{\gamma}$ for the dimension of
    $\chi_{\alpha,\beta,\gamma}$. For every pair of physical indices $i,k$,
    let
    \begin{align}
        D^{\mathrm L}_{ik}
         & =\bigoplus_{\delta,\varepsilon}
        1_{r_{\alpha,\beta}^{\delta}}\otimes
        \bigl(1_{r_{\delta,\gamma}^{\varepsilon}}
        \otimes M_\varepsilon^{ik}\bigr),
        \notag
    \end{align}
    and define $D^{\mathrm R}_{ik}$ by the right-associated multiplicities
    $r_{\beta,\gamma}^{\delta}$ and
    $r_{\alpha,\delta}^{\varepsilon}$. Then
    \begin{align}
        D^{\mathrm L}_{ik}C_{\alpha,\beta,\gamma}
         & =C_{\alpha,\beta,\gamma}D^{\mathrm R}_{ik}.
        \notag
    \end{align}
    This is an identity on the full direct sums. It neither extracts an
    $F$-transformation for a fixed final label nor proves invertibility or
    the pentagon identity.
\end{theorem}
\begin{proof}\leanok
    Length independence makes every positive diagonal factor equal to the
    identity. Substitute the left and right triple-fusion identities into
    (\ref{eq:rfp_triple_comparison}), cancel the two isometries, and apply
    Theorem~\ref{thm:mpdo_operator_product_associativity}.
\end{proof}

\begin{definition}[Common word selectors for the final labels]%
    \label{def:mpdo_final_label_selector_words}
    \lean{MPOTensor.BNTFusionIsometryFamily.HasFinalLabelSelectorWords}
    \leanok
    Fix a common length $S$. The tensors admit common word selectors for
    the final labels at length $S$ if, for each final label $\varepsilon$,
    there are coefficients $c_\varepsilon(w)$, indexed by the words $w$
    of length $S$ in the doubled physical alphabet, such that
    \begin{align}
        \sum_{|w|=S}c_\varepsilon(w)M_\varepsilon^w
         & =1,
        \notag \\
        \sum_{|w|=S}c_\varepsilon(w)M_{\varepsilon'}^w
         & =0
        \notag
    \end{align}
    for every $\varepsilon'\ne\varepsilon$.
    This is the finite-word form of the simultaneous inverse relation
    $B_d^+B_{d'}=\delta_{d,d'}1$ used in the fixed-channel argument of
    \cite{Bultinck2017Anyons}.
\end{definition}

\begin{theorem}[Pair fusion completeness from positive-length selectors]%
    \label{thm:mpdo_pair_fusion_selector_completeness}
    \lean{MPOTensor.BNTFusionIsometryFamily.%
        fusionIsometry_mul_conjTranspose_eq_one_of_lengthIndependent_of_selectorWords}
    \leanok
    \uses{def:mpdo_final_label_selector_words, thm:mpdo_bnt_length_independent_zipper}
    Suppose the structure coefficients
    $c^{(L)}_{\alpha,\beta,\gamma}
        =\tr(\chi_{\alpha,\beta,\gamma}^{L})$, with positive diagonal
    $\chi_{\alpha,\beta,\gamma}$, are independent of the positive chain
    length, and suppose that common selectors exist at a positive length.
    Then every pair fusion isometry is surjective:
    $U_{\alpha,\beta}U_{\alpha,\beta}^{\dagger}=1$.
    The common selectors are the finite-word form of the simultaneous inverse
    used in \cite[line~269]{Bultinck2017Anyons}.
\end{theorem}
\begin{proof}\leanok
    For a word $w$ of the positive selector length, the identity
    $U_{\alpha,\beta}^{\dagger}U_{\alpha,\beta}=1$ gives
    $D^w=U_{\alpha,\beta}(M_\alpha M_\beta)^w
        U_{\alpha,\beta}^{\dagger}$, where
    $D=\bigoplus_\gamma
        1_{r_{\alpha,\beta}^{\gamma}}\otimes M_\gamma$.
    Hence the range projection
    $P=U_{\alpha,\beta}U_{\alpha,\beta}^{\dagger}$ satisfies $PD^w=D^w$.
    Summing the selector polynomials over the final label gives the identity
    on the full direct sum, and therefore $P=1$.
\end{proof}

\begin{theorem}[Pair fusion completeness from a BNT]%
    \label{thm:mpdo_pair_fusion_bnt_completeness}
    \lean{MPOTensor.BNTFusionIsometryFamily.%
        hasFinalLabelSelectorWords_of_hasBlockSelectorWords}
    \lean{MPOTensor.BNTFusionIsometryFamily.%
        fusionIsometry_mul_conjTranspose_eq_one_of_bnt_of_lengthIndependent}
    \leanok
    \uses{thm:cpsv_bnt_blocked_simultaneous_span, thm:mpdo_pair_fusion_selector_completeness}
    Suppose that the tensors $\{M_\alpha\}_\alpha$ with the given fusion
    isometries form a BNT and that the structure coefficients
    $c^{(L)}_{\alpha,\beta,\gamma}
        =\tr(\chi_{\alpha,\beta,\gamma}^{L})$, with positive diagonal
    $\chi_{\alpha,\beta,\gamma}$, are independent of the positive chain
    length. Then every pair fusion isometry is surjective:
    $U_{\alpha,\beta}U_{\alpha,\beta}^{\dagger}=1$.
    In particular, common final-label selectors and the coisometry identity
    follow from these hypotheses. This is a consequence of the common BNT
    word span in \cite[lines~317--345]{Cirac2016MPDO_arXiv}, the fusion
    identity at lines~986--993, and length independence at lines~995--1010.
\end{theorem}
\begin{proof}\leanok
    The basis-of-normal-tensors argument gives a positive common word length
    at which the labelled word evaluations span the full product matrix
    algebra. Choosing the tuple which is the identity in one labelled block
    and zero in every other block gives common final-label selectors. Apply
    Theorem~\ref{thm:mpdo_pair_fusion_selector_completeness}.
\end{proof}
